\chapter{Debug the Build Systematically}
\label{chap:debugging}

\begin{learningoutcomes}
  \item Find the first actionable error.
  \item Reduce a problem to a minimal example.
  \item Distinguish source, package, path, and build-sequence failures.
\end{learningoutcomes}

\section{Use one repair loop}

\begin{enumerate}
  \item Save the source.
  \item Compile the root file.
  \item Read the first error, not the last page of the log.
  \item Inspect the named line and the few lines before it.
  \item Make one repair.
  \item Compile again.
\end{enumerate}

Later errors are often consequences of the first missing brace, command, file,
or environment ending.

\section{Recognise the common messages}

\begin{center}
\begin{tabularx}{\linewidth}{@{}lX@{}}
\toprule
Message fragment & First check \\
\midrule
Undefined control sequence & Command spelling and required package. \\
File not found & Filename, extension, capitalisation, and relative path. \\
Missing \code{\}} inserted & Braces on the current and preceding lines. \\
\code{\textbackslash begin} ended by \code{\textbackslash end} & Environment
names and nesting. \\
Extra alignment tab & Ampersands compared with table columns. \\
Citation undefined & Citation key, bibliography file, and Biber run. \\
Reference undefined & Label spelling and repeated compilation. \\
Overfull box & Long word, URL, equation, table, or code line. \\
\bottomrule
\end{tabularx}
\end{center}

\section{Make a minimal working example}

Copy the failing idea into this file and add material until the error returns:

\begin{latexcode}{Complete example: debugging shell}
\documentclass{article}

% Load only packages needed to reproduce the problem.

\begin{document}
The smallest content that demonstrates the problem.
\end{document}
\end{latexcode}

If the small file works, the problem depends on something removed. If it fails,
the shorter log is easier to interpret.

\section{Compare a broken and repaired line}

\begin{latexcode}{Broken example: reserved characters}
The success rate was 95% in sample_A.
\end{latexcode}

\begin{latexcode}{Corrected line}
The success rate was 95\% in sample\_A.
\end{latexcode}

\begin{latexcode}{Broken example: mismatched environment}
\begin{itemize}
  \item First item
\end{enumerate}
\end{latexcode}

\begin{latexcode}{Corrected environment}
\begin{itemize}
  \item First item
\end{itemize}
\end{latexcode}

\section{Force a clean rebuild}

\begin{terminalcode}{Clean and rebuild}
latexmk -C main.tex
latexmk -pdf main.tex
\end{terminalcode}

Clean only generated files. Never use a broad deletion command to diagnose a
build.

\section{Check for silent layout failures}

A successful compilation can still contain clipped tables, tiny figures,
empty pages, or unresolved warnings. Search the log and inspect the PDF:

\begin{terminalcode}{Search a build log}
rg "Overfull|undefined|multiply defined" main.log
\end{terminalcode}

\begin{codelab}
Introduce one missing brace, one wrong environment name, and one bad path into
separate copies of a working file. Repair each using only the first error and
the nearby source.
\end{codelab}

\begin{chapterchecklist}
  \item The first error was repaired before later messages.
  \item Only one change was tested at a time.
  \item A minimal example was used when the cause was unclear.
  \item The PDF was inspected even after a successful build.
\end{chapterchecklist}
